% !TEX program = pdflatex
% THEME 2 — Facile — Exercices
\documentclass[11pt,a4paper]{article}
\usepackage[utf8]{inputenc}
\usepackage[T1]{fontenc}
\usepackage{lmodern}
\usepackage[french]{babel}
\usepackage{amsmath,amssymb,mathtools}
\usepackage{icomma}
\usepackage{siunitx}
\sisetup{output-decimal-marker={,},per-mode=symbol}
\DeclareSIUnit{\atm}{atm}
\DeclareSIUnit{\bar}{bar}
\DeclareSIUnit{\mmHg}{mmHg}
\usepackage[version=4]{mhchem}
\usepackage{xcolor}
\usepackage{colortbl}
\usepackage[most]{tcolorbox}
\usepackage{enumitem}
\usepackage{array}
\usepackage{booktabs}
\usepackage{fancyhdr}
\usepackage[a4paper,margin=2.1cm,headheight=26pt,headsep=10pt]{geometry}

\definecolor{primary}{RGB}{150,98,162}
\definecolor{primarydark}{RGB}{92,52,110}

\newtcolorbox{enonce}[1][]{enhanced,breakable,colback=primary!4,colframe=primary,
  boxrule=0.8pt,arc=2pt,left=3mm,right=3mm,top=2mm,bottom=2mm,
  fonttitle=\bfseries,coltitle=white,title={#1}}

\pagestyle{fancy}\fancyhf{}
\fancyhead[L]{\small\textcolor{primarydark}{\textbf{Fiche 7} — Loi des gaz}}
\fancyhead[R]{\small\textcolor{primarydark}{\textsc{Thème 2 · Facile · Exercices}}}
\fancyfoot[C]{\small\thepage}
\renewcommand{\headrulewidth}{0.8pt}
\renewcommand{\headrule}{\hbox to\headwidth{\color{primary}\leaders\hrule height \headrulewidth\hfill}}
\setlength{\parindent}{0pt}\setlength{\parskip}{3pt}

\newcounter{exo}
\newenvironment{exo}[1][]{\refstepcounter{exo}\medskip
  \noindent{\color{primarydark}\bfseries Exercice \theexo.}\ \textit{#1}\par\nopagebreak\smallskip}{\par}

\begin{document}

\begin{center}
{\Large\bfseries\color{primarydark} Thème 2 — Niveau Facile}\\[2pt]
{\color{primary}\itshape Isoler chaque variable de $pV=nRT$ en une étape}
\end{center}
\textcolor{primary}{\hrule height 1pt}
\medskip

\begin{enonce}[Rappels utiles]
$pV=nRT$ avec $T$ en \textbf{kelvins} ($T=\theta+273$). \\
Constante : $R=\qty{0.082}{\atm\litre\per\mole\per\kelvin}$ (avec $p$ en \si{\atm}, $V$ en \si{\litre}) ;
ou $R=\qty{8.314}{\kilo\pascal\litre\per\mole\per\kelvin}$ (avec $p$ en \si{\kilo\pascal}, $V$ en \si{\litre}).
\end{enonce}

\begin{exo}[Combien de moles ?]
Un récipient de volume $V=\qty{10.0}{\litre}$ contient un gaz sous une pression $p=\qty{2.0}{\atm}$ à la
température $T=\qty{300}{\kelvin}$. Détermine la quantité de matière $n$ de gaz présente.
\end{exo}

\begin{exo}[Quelle pression ?]
On enferme $n=\qty{0.50}{\mole}$ de dioxygène \ce{O2} dans une bouteille rigide de volume
$V=\qty{5.0}{\litre}$ portée à $\theta=\qty{27}{\degreeCelsius}$. Quelle pression $p$ règne à l'intérieur ?
\end{exo}

\begin{exo}[Quel volume ?]
Quel volume $V$ occupent $n=\qty{2.0}{\mole}$ de gaz sous une pression $p=\qty{1.0}{\atm}$ à une
température $T=\qty{273}{\kelvin}$ ?
\end{exo}

\begin{exo}[Quelle température ?]
Un gaz constitué de $n=\qty{1.0}{\mole}$ occupe un volume $V=\qty{22.4}{\litre}$ sous une pression
$p=\qty{1.0}{\atm}$. Quelle est sa température $T$ (en kelvins, puis en degrés Celsius) ?
\end{exo}

\begin{exo}[Choisir le bon $R$]
Un gaz occupe $V=\qty{2.0}{\litre}$ sous une pression $p=\qty{101.3}{\kilo\pascal}$ à
$T=\qty{273}{\kelvin}$.
\begin{enumerate}[label=(\alph*),leftmargin=2em,itemsep=1pt]
  \item La pression est ici en \si{\kilo\pascal} : quelle valeur de $R$ dois-tu utiliser ?
  \item Calcule la quantité de matière $n$.
\end{enumerate}
\end{exo}

\end{document}
